%-------------------------------------------------------------------------------
% §7 CAHIERS DES CHARGES, SPÉCIFICATIONS ET ARCHITECTURES DE SI
%-------------------------------------------------------------------------------
\newpage
\cvsection{7. Cahiers des charges, spécifications et architectures SI}

\cvgroupreset

\begin{cvgroup}
\cvsubsection{Architecture ASYCUDA – EBMS – POS-OBR}

\begin{cvmetadata}
\textbf{Client :} Office Burundais des Recettes (OBR)
\end{cvmetadata}

\begin{cvprojectsummary}
Production des livrables d'architecture et de spécification pour l'intégration du système douanier ASYCUDA, de la plateforme fiscale EBMS et de la solution POS-OBR, dans le cadre de la digitalisation des services fiscaux publics au Burundi. Modélisation UML par diagrammes de séquences des flux inter-systèmes et des interactions entre acteurs métier.
\end{cvprojectsummary}

\begin{cvdetailslist}{Contributions principales}
  \item {Rédaction du cahier des charges fonctionnel et technique de l'intégration ASYCUDA–EBMS–POS-OBR.}
  \item {Modélisation UML (diagrammes de séquences) et diagrammes de flux inter-systèmes.}
  \item {Rédaction des spécifications techniques, contrats d'API JSON et modèles de données PostgreSQL.}
  \item {Définition du cadre de gouvernance des données publiques, de validation et de traçabilité.}
  \item {Élaboration de la feuille de route de déploiement et du plan de mise en œuvre sur GCP/Docker.}
  \item {Analyse d'impact et vérification de la conformité aux exigences fonctionnelles et non-fonctionnelles.}
\end{cvdetailslist}

\cvdetailblock{Technologies}{UML (diagrammes de séquences) · Dart · Flutter · Docker · Google Cloud Platform · PostgreSQL}

\cvdetailblock{Compétences démontrées}{Architecture de SI · Cahier des charges · Spécifications fonctionnelles et techniques · Modélisation UML · Gouvernance des données publiques · Institution publique · E-gouvernement · Feuille de route}
\end{cvgroup}

\begin{cvgroup}
\cvsubsection{Architecture CMMS Nemidis}

\begin{cvprojectsummary}
Conception de l'architecture globale d'une plateforme multi-institutions de gestion de maintenance des Machines de Facturation Électronique (MFE), avec modélisation des workflows et des interactions entre acteurs.
\end{cvprojectsummary}

\begin{cvdetailslist}{Contributions principales}
  \item {Rédaction du cahier des charges fonctionnel et technique.}
  \item {Définition de l'architecture multi-tenant, du modèle de données PostgreSQL et des contrats d'API.}
  \item {Conception des workflows de maintenance, diagnostic, SLA et facturation.}
  \item {Définition des règles métier, contrôles d'accès et mécanismes de traçabilité.}
  \item {Planification des développements, supervision de l'implémentation et déploiement sur GCP/Docker.}
\end{cvdetailslist}

\cvdetailblock{Technologies}{UML (diagrammes de séquences) · Dart · Flutter · Docker · Google Cloud Platform · PostgreSQL}

\cvdetailblock{Compétences démontrées}{Architecture de SI · Cahier des charges · Modélisation de données · Gouvernance des données · Supervision de projet · Déploiement cloud}
\end{cvgroup}
